\subsection{商空间的庞加莱群胚}

设 X 是赋予离散群 G 连续作用的拓扑空间；令 Y = X/G，并将 \(f: X \to Y\) 记为规范映射。记 \(\lvert G\rvert\) 为群胚 \(\varpi(G)\)；其顶点集为 G；对于 \(g, g' \in G\)，存在从 g 到 \(g'\) 的唯一箭，当且仅当 \(g' = g\)，否则不存在。通过取基本群胚，作用 \(m: G \times X \to X\) 诱导出群胚态射 \(\varpi(m): \lvert G\rvert \times \varpi(X) \to \varpi(X)\)。令 \(\varpi(X)/G\) 为群胚态射 \(\varpi(m)\) 和 \(\varpi(\mathrm{pr}_{2})\) 的余均衡子，这两个态射从 \(\lvert G\rvert \times \varpi(X)\) 到 \(\varpi(X)\)；记 \(\beta: \varpi(X) \to \varpi(X)/G\) 为规范群胚态射。我们有 \(f \circ m = f \circ pr_{2}\)，因此 \(\varpi(f) \circ \varpi(m) = \varpi(f) \circ \varpi(\mathrm{pr}_{2})\)。根据余均衡子的泛性质，存在唯一群胚态射 \(\varpi''(f): \varpi(X)/G \to \varpi(Y)\)，使得 \(\varpi(f) = \varpi''(f) \circ \beta\)。

定理 3. --- 设 X 是可解圈拓扑空间，G 是适当地作用于 X 的离散群；设 \(f\colon X \to X/G\) 为规范满射。上面引入的规范群胚态射 \(\varpi''(f)\colon \varpi(X)/G \to \varpi(X/G)\) 是同构。

记 \(\operatorname{Coeg}(f)\) 为从 \(\varpi(\mathrm{X} \times_{\mathrm{Y}} \mathrm{X})\) 到 \(\varpi(\mathrm{X})\) 的两个群胚态射的余均衡子，这两个态射由投影 \(pr_{1}\) 和 \(pr_{2}\) 诱导；令 \(\gamma: \varpi(\mathrm{X}) \to \operatorname{Coeg}(f)\) 为规范群胚态射。由于映射 \((m, pr_{2})\) : \(G \times X \to X \times X\) 的像是 \(X \times_{Y} X\) 在 \(X \times X\) 中的子空间，群胚态射 \(\gamma \circ \varpi(m)\) 和 \(\gamma \circ \varpi(pr_{2})\)（从 \(\lvert G\rvert \times \varpi(X)\) 到 \(\operatorname{Coeg}(f)\)）相等。根据 \(\varpi(\mathrm{X})/\mathrm{G}\) 的泛性质，存在唯一群胚态射 \(\alpha: \varpi(\mathrm{X})/\mathrm{G} \to \operatorname{Coeg}(f)\)，使得 \(\gamma = \alpha \circ \beta\)。

我们也将唯一的群胚态射 \(\varpi'(f)\) 记为从 \(\operatorname{Coeg}(f)\) 到 \(\varpi(Y)\) 且满足 \(\varpi(f) = \varpi'(f) \circ \gamma\) 的态射。根据 IV, 第~399~页，例 2，IV, 第~398~页定理 1 的假设得到满足，且态射 \(\varpi'(f)\) 是同构。由于

\[
\varpi (f) = \varpi^ {\prime} (f) \circ \gamma = \varpi^ {\prime} (f) \circ \alpha \circ \beta = \varpi^ {\prime \prime} (f) \circ \beta ,
\]

有 \(\varpi''(f) = \varpi'(f) \circ \alpha\)。因此，要证明定理 3，只需证明态射 \(\alpha\) 是同构。

引理 4. --- 对每条道路 \(c = (c_1, c_2)\)，若其属于 \(X \times_Y X\)，则有 \(\beta([c_1]) = \beta([c_2])\)。

设 c 是这样的一条道路。

设 \(x \in X\)，令 \(K_x\) 为 x 在 G 中的稳定子。根据 TG, III, 第~32~页命题 8，存在一个开邻域 \(U_x\)，它是 \(x\) 在 X 中的邻域，使得 \(K_x \cdot U_x = U_x\)，

\(g \cdot U_x \cap U_x = \varnothing\) 对所有 \(g \in G - K_x\) 成立，并且映射 f 诱导出从 \(U_x/K_x\) 到 Y 中开邻域 \(V_x\) 的同胚，其中包含 \(f(x)\)。由于 X 是可解圈的，且映射 f 在 \(U_x\) 上的限制是开且闭的，不妨进一步假设 \(U_x\) 是连通的，并且规范同态 \(\pi_1(U_x, x) \to \pi_1(X, x)\) 的像化为单位元。这样构造的开集 \((V_x)_{x \in X}\) 构成 Y 的一个开覆盖。根据 III, 第~272~页引理 4，将其应用于紧空间 I 以及开集 \((f \circ c_1)^{-1}(V_x)\)（\(x \in X\)），存在整数 \(n \geqslant 1\)，使得对每个 \(i \in \{1, \ldots, n\}\)，存在 X 中一点 \(x_i\)，使得 \(c_1([i-1/n, i/n])\) 包含于 \(f(V_{x_i})\)。由于 \(f \circ c_1 = f \circ c_2\)，\(c_2([i-1/n, i/n])\) 也包含于 \(f(V_{x_i})\)。

对于 \(j = 1\) 或 2 以及 \(i \in \{1, \ldots, n\}\)，记 \(c_{j,i}\) 为 X 中由 \(t \mapsto c_j(\frac{i+t-1}{n})\) 定义的道路；有 \(c_j = c_{j,1} * \cdots * c_{j,n}\)（III, 第~291~页，注 1）。因此，要证明 \(\beta([c_1]) = \beta([c_2])\)，只需证明 \(\beta([c_{1,i}]) = \beta([c_{2,i}])\)，对每个整数 \(i \in \{1, \ldots, n\}\) 都如此。将道路 \((c_1, c_2)\) 替换为道路 \((c_{1,i}, c_{2,i})\)，于是可假设存在 \(x \in X\)，使得 \((f \circ c_1)([0,1]) \subset V_x\)。

\(V_{x}\) 在 f 下的逆像是连通分支 \(g \cdot U_{x}\) 的不交并，其中 g 遍历 G 中 \(G/K_{x}\) 的一组代表。对于 i = 1 或 2，取 G 中元素 \(g_{i}\)，使得点 \(x_{i} = g_{i} \cdot c_{i}(0)\) 属于 \(U_{x}\)。于是道路 \(g_{i} \cdot c_{i}\) 的像包含于 \(U_{x}\)。根据群胚 \(\varpi(\mathrm{X})/\mathrm{G}\) 的定义，有 \(\beta([g_{i} \cdot c_{i}]) = \beta([c_{i}])\)，从而可以假设道路 \(c_{1}\) 和 \(c_{2}\) 的像包含于 \(U_{x}\)，且 \(g_{1} = g_{2} = e\)。

对于 \(s=0\) 或 1，令 \(d_{s}\) 为 \(\mathrm{U}_{x}\) 中连接 \(x\) 与 \(c_{1}(s)\) 的道路，并令 \(g_{s}\) 为 \(\mathrm{K}_{x}\) 中满足 \(g_{s}\cdot c_{2}(s)=c_{1}(s)\) 的元素。道路 \(d_{0}*c_{1}*\overline{d_{1}}\) 和 \(d_{0}*(g_{0}\cdot c_{2})*(g_{0}g_{1}^{-1}\cdot\overline{d_{1}})\) 是以 \(x\) 为基点、位于 \(\mathrm{U}_{x}\) 中的圈。因此，它们在 X 中都严格同伦于以 \(x\) 为基点的常值圈，因为规范同态 \(\pi_{1}(\mathrm{U}_{x},x)\to\pi_{1}(\mathrm{X},x)\) 的像化为单位元。尤其，它们的类在群胚态射 \(\beta\) 下有相同的像，于是

\[
\beta ([ d _ {0} ]) \beta ([ c _ {1} ]) \beta ([ d _ {1} ]) ^ {- 1} = \beta ([ d _ {0} ]) \beta ([ g _ {0} \cdot c _ {2} ]) \beta ([ g _ {0} g _ {1} ^ {- 1} \cdot d _ {1} ]) ^ {- 1}.
\]

由于 \(\beta([g \cdot c]) = \beta([c])\) 对每个元素 \(g \in G\) 以及每条道路 \(c\)，它位于 \(X\) 中，都成立，故有 \(\beta([c_1]) = \beta([c_2])\)，引理得证。

根据引理，两个群胚态射 \(\beta \circ \varpi(\mathrm{pr}_{1})\) 和 \(\beta \circ \varpi(\mathrm{pr}_{2})\)，它们从 \(\varpi(\mathrm{X} \times_{\mathrm{Y}} \mathrm{X})\) 到 \(\operatorname{Coeg}(f)\)，相等。根据余均衡子的泛性质，因此存在唯一群胚态射 \(\alpha'\)：\(\operatorname{Coeg}(f) \to \varpi(\mathrm{X})/\mathrm{G}\)，使得 \(\beta = \alpha' \circ \gamma\)。态射 \(\alpha' \circ \alpha\) 是唯一的群胚态射 \(\varphi\)，它作用于 \(\varpi(\mathrm{X})/\mathrm{G}\) 并映到自身，且满足 \(\varphi \circ \beta = \beta\)；因此有 \(\alpha' \circ \alpha = \operatorname{Id}_{\varpi(\mathrm{X})/\mathrm{G}}\)。同样，\(\alpha \circ \alpha' = \operatorname{Id}_{\operatorname{Coeg}(f)}\)。所以，\(\alpha\) 是同构。

